%% key-resolution-cost.tex
%%
%% "It Is Not the Cryptography: A Discarded Exception Dominates Token
%%  Validation Cost in a Node.js Service"
%%
%% Supersedes not-the-cryptography.tex, whose mechanism claim, Table 2
%% explanation and stated runtime version are corrected here.
%%
%% EVERY numeric value in this document is a macro from keypath_macros.tex,
%% generated by analysis/make_keypath_macros.py from files under data/runs/ and
%% survey/. No number is typed inline. keypath_macros_provenance.csv maps each
%% macro to the file it was read from. A quantity that was not measured is a
%% \PLACEHOLDER and typesets red.
\documentclass[pdflatex,sn-basic,twocolumn]{sn-jnl}

\usepackage{lmodern}
\usepackage{manyfoot}
\usepackage{amsmath}
\usepackage{amssymb}
\usepackage{graphicx}
\usepackage{booktabs}
\usepackage{url}
\usepackage{xcolor}

% keypath_macros.tex -- GENERATED. Do not edit.
% Regenerate: .venv/bin/python -m analysis.make_keypath_macros
% Every value below is read from a file under data/runs/ or survey/.
\providecommand{\PLACEHOLDER}[1]{\textcolor{red}{\textbf{[MISSING: #1]}}}

\newcommand{\AdjudicatedFalsePositive}{5}% survey/data/hand_adjudication.csv
\newcommand{\AdjudicatedGenuine}{7}% survey/data/hand_adjudication.csv
\newcommand{\AdjudicatedGenuineProjects}{6}% survey/data/hand_adjudication.csv
\newcommand{\AdjudicatedTotal}{13}% survey/data/hand_adjudication.csv
\newcommand{\AdjudicatedUnresolved}{1}% survey/data/hand_adjudication.csv
\newcommand{\ContainerVersusHostRatio}{0.92}% data/runs/2026-09-17T08-57-13Z-keypath-runtime-matrix/keypath_stats.json
\newcommand{\CounterFilesFetched}{191}% survey/data/counterexamples.json
\newcommand{\CounterWithVerify}{120}% survey/data/counterexamples.json
\newcommand{\DecompExplainedPct}{77.9}% data/runs/2026-09-17T08-42-17Z-keypath-mechanism/keypath_stats.json
\newcommand{\DecompObserved}{23.92}% data/runs/2026-09-17T08-42-17Z-keypath-mechanism/keypath_stats.json
\newcommand{\DecompPredicted}{18.62}% data/runs/2026-09-17T08-42-17Z-keypath-mechanism/keypath_stats.json
\newcommand{\DecompResidual}{5.29}% data/runs/2026-09-17T08-42-17Z-keypath-mechanism/keypath_stats.json
\newcommand{\DecompResidualCiHi}{5.50}% data/runs/2026-09-17T08-42-17Z-keypath-mechanism/keypath_stats.json
\newcommand{\DecompResidualCiLo}{4.92}% data/runs/2026-09-17T08-42-17Z-keypath-mechanism/keypath_stats.json
\newcommand{\HostCallsPerInvocation}{40000}% data/runs/2026-09-17T08-42-17Z-keypath-mechanism/keypath_stats.json
\newcommand{\HostCreateSecretKey}{0.583}% data/runs/2026-09-17T08-42-17Z-keypath-mechanism/keypath_stats.json
\newcommand{\HostCreateSecretKeyCiHi}{0.58}% data/runs/2026-09-17T08-42-17Z-keypath-mechanism/keypath_stats.json
\newcommand{\HostCreateSecretKeyCiLo}{0.54}% data/runs/2026-09-17T08-42-17Z-keypath-mechanism/keypath_stats.json
\newcommand{\HostCvMax}{3.5}% data/runs/2026-09-17T08-42-17Z-keypath-mechanism/keypath_stats.json
\newcommand{\HostDecodeOnly}{1.29}% data/runs/2026-09-17T08-42-17Z-keypath-mechanism/keypath_stats.json
\newcommand{\HostDecodeOnlyCiHi}{1.29}% data/runs/2026-09-17T08-42-17Z-keypath-mechanism/keypath_stats.json
\newcommand{\HostDecodeOnlyCiLo}{1.25}% data/runs/2026-09-17T08-42-17Z-keypath-mechanism/keypath_stats.json
\newcommand{\HostHmacKeyObject}{2.17}% data/runs/2026-09-17T08-42-17Z-keypath-mechanism/keypath_stats.json
\newcommand{\HostHmacKeyObjectCiHi}{2.17}% data/runs/2026-09-17T08-42-17Z-keypath-mechanism/keypath_stats.json
\newcommand{\HostHmacKeyObjectCiLo}{2.12}% data/runs/2026-09-17T08-42-17Z-keypath-mechanism/keypath_stats.json
\newcommand{\HostHmacString}{2.21}% data/runs/2026-09-17T08-42-17Z-keypath-mechanism/keypath_stats.json
\newcommand{\HostHmacStringCiHi}{2.25}% data/runs/2026-09-17T08-42-17Z-keypath-mechanism/keypath_stats.json
\newcommand{\HostHmacStringCiLo}{2.21}% data/runs/2026-09-17T08-42-17Z-keypath-mechanism/keypath_stats.json
\newcommand{\HostInvocations}{15}% data/runs/2026-09-17T08-42-17Z-keypath-mechanism/keypath_stats.json
\newcommand{\HostJwtPreparsed}{6.04}% data/runs/2026-09-17T08-42-17Z-keypath-mechanism/keypath_stats.json
\newcommand{\HostJwtPreparsedCiHi}{6.12}% data/runs/2026-09-17T08-42-17Z-keypath-mechanism/keypath_stats.json
\newcommand{\HostJwtPreparsedCiLo}{5.96}% data/runs/2026-09-17T08-42-17Z-keypath-mechanism/keypath_stats.json
\newcommand{\HostJwtRsPemString}{52.58}% data/runs/2026-09-17T08-42-17Z-keypath-mechanism/keypath_stats.json
\newcommand{\HostJwtRsPemStringCiHi}{52.71}% data/runs/2026-09-17T08-42-17Z-keypath-mechanism/keypath_stats.json
\newcommand{\HostJwtRsPemStringCiLo}{52.42}% data/runs/2026-09-17T08-42-17Z-keypath-mechanism/keypath_stats.json
\newcommand{\HostJwtRsPreparsed}{25.33}% data/runs/2026-09-17T08-42-17Z-keypath-mechanism/keypath_stats.json
\newcommand{\HostJwtRsPreparsedCiHi}{25.50}% data/runs/2026-09-17T08-42-17Z-keypath-mechanism/keypath_stats.json
\newcommand{\HostJwtRsPreparsedCiLo}{25.29}% data/runs/2026-09-17T08-42-17Z-keypath-mechanism/keypath_stats.json
\newcommand{\HostJwtString}{29.92}% data/runs/2026-09-17T08-42-17Z-keypath-mechanism/keypath_stats.json
\newcommand{\HostJwtStringCiHi}{30.12}% data/runs/2026-09-17T08-42-17Z-keypath-mechanism/keypath_stats.json
\newcommand{\HostJwtStringCiLo}{29.75}% data/runs/2026-09-17T08-42-17Z-keypath-mechanism/keypath_stats.json
\newcommand{\HostJwtStringSafe}{7.21}% data/runs/2026-09-17T08-42-17Z-keypath-mechanism/keypath_stats.json
\newcommand{\HostJwtStringSafeCiHi}{7.33}% data/runs/2026-09-17T08-42-17Z-keypath-mechanism/keypath_stats.json
\newcommand{\HostJwtStringSafeCiLo}{7.12}% data/runs/2026-09-17T08-42-17Z-keypath-mechanism/keypath_stats.json
\newcommand{\HostNodeVersion}{26.6.0}% data/runs/2026-09-17T08-42-17Z-keypath-mechanism/keypath_stats.json
\newcommand{\HostProbeShareOfCall}{60}% data/runs/2026-09-17T08-42-17Z-keypath-mechanism/keypath_stats.json
\newcommand{\HostProbeSucceeds}{13.83}% data/runs/2026-09-17T08-42-17Z-keypath-mechanism/keypath_stats.json
\newcommand{\HostProbeSucceedsCiHi}{13.92}% data/runs/2026-09-17T08-42-17Z-keypath-mechanism/keypath_stats.json
\newcommand{\HostProbeSucceedsCiLo}{13.62}% data/runs/2026-09-17T08-42-17Z-keypath-mechanism/keypath_stats.json
\newcommand{\HostProbeThrows}{18.04}% data/runs/2026-09-17T08-42-17Z-keypath-mechanism/keypath_stats.json
\newcommand{\HostProbeThrowsCiHi}{18.25}% data/runs/2026-09-17T08-42-17Z-keypath-mechanism/keypath_stats.json
\newcommand{\HostProbeThrowsCiLo}{17.83}% data/runs/2026-09-17T08-42-17Z-keypath-mechanism/keypath_stats.json
\newcommand{\HostRsStringPenalty}{27.29}% data/runs/2026-09-17T08-42-17Z-keypath-mechanism/keypath_stats.json
\newcommand{\HostRsStringPenaltyCiHi}{27.37}% data/runs/2026-09-17T08-42-17Z-keypath-mechanism/keypath_stats.json
\newcommand{\HostRsStringPenaltyCiLo}{27.00}% data/runs/2026-09-17T08-42-17Z-keypath-mechanism/keypath_stats.json
\newcommand{\HostRsStringPenaltyRounds}{15}% data/runs/2026-09-17T08-42-17Z-keypath-mechanism/keypath_stats.json
\newcommand{\HostSafePatchRatio}{4.15}% data/runs/2026-09-17T08-42-17Z-keypath-mechanism/keypath_stats.json
\newcommand{\HostSafePatchResidual}{1.21}% data/runs/2026-09-17T08-42-17Z-keypath-mechanism/keypath_stats.json
\newcommand{\HostSafePatchResidualCiHi}{1.25}% data/runs/2026-09-17T08-42-17Z-keypath-mechanism/keypath_stats.json
\newcommand{\HostSafePatchResidualCiLo}{1.08}% data/runs/2026-09-17T08-42-17Z-keypath-mechanism/keypath_stats.json
\newcommand{\HostSafePatchResidualRounds}{15}% data/runs/2026-09-17T08-42-17Z-keypath-mechanism/keypath_stats.json
\newcommand{\HostSafePatchSaving}{22.71}% data/runs/2026-09-17T08-42-17Z-keypath-mechanism/keypath_stats.json
\newcommand{\HostSafePatchSavingCiHi}{22.88}% data/runs/2026-09-17T08-42-17Z-keypath-mechanism/keypath_stats.json
\newcommand{\HostSafePatchSavingCiLo}{22.54}% data/runs/2026-09-17T08-42-17Z-keypath-mechanism/keypath_stats.json
\newcommand{\HostSafePatchSavingRounds}{15}% data/runs/2026-09-17T08-42-17Z-keypath-mechanism/keypath_stats.json
\newcommand{\HostStringPenalty}{23.92}% data/runs/2026-09-17T08-42-17Z-keypath-mechanism/keypath_stats.json
\newcommand{\HostStringPenaltyCiHi}{24.00}% data/runs/2026-09-17T08-42-17Z-keypath-mechanism/keypath_stats.json
\newcommand{\HostStringPenaltyCiLo}{23.58}% data/runs/2026-09-17T08-42-17Z-keypath-mechanism/keypath_stats.json
\newcommand{\HostStringPenaltyRounds}{15}% data/runs/2026-09-17T08-42-17Z-keypath-mechanism/keypath_stats.json
\newcommand{\HostTimerOverhead}{0.042}% data/runs/2026-09-17T08-42-17Z-keypath-mechanism/keypath_stats.json
\newcommand{\HostTimerOverheadCiHi}{0.04}% data/runs/2026-09-17T08-42-17Z-keypath-mechanism/keypath_stats.json
\newcommand{\HostTimerOverheadCiLo}{0.04}% data/runs/2026-09-17T08-42-17Z-keypath-mechanism/keypath_stats.json
\newcommand{\IdentityMismatches}{0}% data/runs/2026-09-17T08-57-13Z-keypath-runtime-matrix/environments.csv
\newcommand{\ManifestCounterFiles}{191}% survey/data/corpus_manifest.csv
\newcommand{\ManifestFiles}{502}% survey/data/corpus_manifest.csv
\newcommand{\ManifestHeadResolved}{501}% survey/data/corpus_manifest.csv
\newcommand{\ManifestRepos}{457}% survey/data/corpus_manifest.csv
\newcommand{\ManifestSampleFiles}{311}% survey/data/corpus_manifest.csv
\newcommand{\MatrixCallsPerInvocation}{20000}% data/runs/2026-09-17T08-57-13Z-keypath-runtime-matrix/keypath_stats.json
\newcommand{\MatrixCvMax}{20.2}% data/runs/2026-09-17T08-57-13Z-keypath-runtime-matrix/keypath_stats.json
\newcommand{\MatrixCvMaxCondition}{probe\_succeeds}% data/runs/2026-09-17T08-57-13Z-keypath-runtime-matrix/keypath_stats.json
\newcommand{\MatrixCvMaxEnvironment}{node:26.6.0-bookworm}% data/runs/2026-09-17T08-57-13Z-keypath-runtime-matrix/keypath_stats.json
\newcommand{\MatrixCvRatio}{5.7}% data/runs/2026-09-17T08-57-13Z-keypath-runtime-matrix/keypath_stats.json
\newcommand{\MatrixInvocations}{10}% data/runs/2026-09-17T08-57-13Z-keypath-runtime-matrix/keypath_stats.json
\newcommand{\NodeEighteenDecodeOnly}{1.62}% data/runs/2026-09-17T08-57-13Z-keypath-runtime-matrix/keypath_stats.json
\newcommand{\NodeEighteenHmacString}{2.42}% data/runs/2026-09-17T08-57-13Z-keypath-runtime-matrix/keypath_stats.json
\newcommand{\NodeEighteenImageId}{8d6421d663b4}% data/runs/runtime_identity.csv
\newcommand{\NodeEighteenJwtPreparsed}{7.58}% data/runs/2026-09-17T08-57-13Z-keypath-runtime-matrix/keypath_stats.json
\newcommand{\NodeEighteenJwtString}{394.71}% data/runs/2026-09-17T08-57-13Z-keypath-runtime-matrix/keypath_stats.json
\newcommand{\NodeEighteenNodeLiteral}{18.20.8}% data/runs/2026-09-17T08-57-13Z-keypath-runtime-matrix/keypath_stats.json
\newcommand{\NodeEighteenNodeMajor}{18}% data/runs/2026-09-17T08-57-13Z-keypath-runtime-matrix/keypath_stats.json
\newcommand{\NodeEighteenOpensslLiteral}{3.0.16}% data/runs/2026-09-17T08-57-13Z-keypath-runtime-matrix/keypath_stats.json
\newcommand{\NodeEighteenOpensslSeries}{3.0}% data/runs/2026-09-17T08-57-13Z-keypath-runtime-matrix/keypath_stats.json
\newcommand{\NodeEighteenProbeSucceeds}{115.83}% data/runs/2026-09-17T08-57-13Z-keypath-runtime-matrix/keypath_stats.json
\newcommand{\NodeEighteenProbeThrows}{373.79}% data/runs/2026-09-17T08-57-13Z-keypath-runtime-matrix/keypath_stats.json
\newcommand{\NodeEighteenProbeThrowsCiHi}{375.79}% data/runs/2026-09-17T08-57-13Z-keypath-runtime-matrix/keypath_stats.json
\newcommand{\NodeEighteenProbeThrowsCiLo}{371.50}% data/runs/2026-09-17T08-57-13Z-keypath-runtime-matrix/keypath_stats.json
\newcommand{\NodeEighteenVeight}{10}% data/runs/runtime_identity.csv
\newcommand{\NodeEighteenVeightFull}{10.2.154.26-node.39}% data/runs/runtime_identity.csv
\newcommand{\NodeTwentyDecodeOnly}{2.19}% data/runs/2026-09-17T08-57-13Z-keypath-runtime-matrix/keypath_stats.json
\newcommand{\NodeTwentyHmacString}{2.33}% data/runs/2026-09-17T08-57-13Z-keypath-runtime-matrix/keypath_stats.json
\newcommand{\NodeTwentyImageId}{fb4cd12c85ee}% data/runs/runtime_identity.csv
\newcommand{\NodeTwentyJwtPreparsed}{7.25}% data/runs/2026-09-17T08-57-13Z-keypath-runtime-matrix/keypath_stats.json
\newcommand{\NodeTwentyJwtString}{437.40}% data/runs/2026-09-17T08-57-13Z-keypath-runtime-matrix/keypath_stats.json
\newcommand{\NodeTwentyNodeLiteral}{20.20.2}% data/runs/2026-09-17T08-57-13Z-keypath-runtime-matrix/keypath_stats.json
\newcommand{\NodeTwentyNodeMajor}{20}% data/runs/2026-09-17T08-57-13Z-keypath-runtime-matrix/keypath_stats.json
\newcommand{\NodeTwentyOpensslLiteral}{3.0.19}% data/runs/2026-09-17T08-57-13Z-keypath-runtime-matrix/keypath_stats.json
\newcommand{\NodeTwentyOpensslSeries}{3.0}% data/runs/2026-09-17T08-57-13Z-keypath-runtime-matrix/keypath_stats.json
\newcommand{\NodeTwentyProbeSucceeds}{142.38}% data/runs/2026-09-17T08-57-13Z-keypath-runtime-matrix/keypath_stats.json
\newcommand{\NodeTwentyProbeThrows}{420.62}% data/runs/2026-09-17T08-57-13Z-keypath-runtime-matrix/keypath_stats.json
\newcommand{\NodeTwentyProbeThrowsCiHi}{422.56}% data/runs/2026-09-17T08-57-13Z-keypath-runtime-matrix/keypath_stats.json
\newcommand{\NodeTwentyProbeThrowsCiLo}{419.52}% data/runs/2026-09-17T08-57-13Z-keypath-runtime-matrix/keypath_stats.json
\newcommand{\NodeTwentySixGlibcDecodeOnly}{1.67}% data/runs/2026-09-17T08-57-13Z-keypath-runtime-matrix/keypath_stats.json
\newcommand{\NodeTwentySixGlibcHmacString}{3.19}% data/runs/2026-09-17T08-57-13Z-keypath-runtime-matrix/keypath_stats.json
\newcommand{\NodeTwentySixGlibcImageId}{6e0a69f43d2b}% data/runs/runtime_identity.csv
\newcommand{\NodeTwentySixGlibcJwtPreparsed}{6.04}% data/runs/2026-09-17T08-57-13Z-keypath-runtime-matrix/keypath_stats.json
\newcommand{\NodeTwentySixGlibcJwtString}{27.96}% data/runs/2026-09-17T08-57-13Z-keypath-runtime-matrix/keypath_stats.json
\newcommand{\NodeTwentySixGlibcNodeLiteral}{26.6.0}% data/runs/2026-09-17T08-57-13Z-keypath-runtime-matrix/keypath_stats.json
\newcommand{\NodeTwentySixGlibcNodeMajor}{26}% data/runs/2026-09-17T08-57-13Z-keypath-runtime-matrix/keypath_stats.json
\newcommand{\NodeTwentySixGlibcOpensslLiteral}{3.5.7}% data/runs/2026-09-17T08-57-13Z-keypath-runtime-matrix/keypath_stats.json
\newcommand{\NodeTwentySixGlibcOpensslSeries}{3.5}% data/runs/2026-09-17T08-57-13Z-keypath-runtime-matrix/keypath_stats.json
\newcommand{\NodeTwentySixGlibcProbeSucceeds}{12.12}% data/runs/2026-09-17T08-57-13Z-keypath-runtime-matrix/keypath_stats.json
\newcommand{\NodeTwentySixGlibcProbeThrows}{16.58}% data/runs/2026-09-17T08-57-13Z-keypath-runtime-matrix/keypath_stats.json
\newcommand{\NodeTwentySixGlibcProbeThrowsCiHi}{17.88}% data/runs/2026-09-17T08-57-13Z-keypath-runtime-matrix/keypath_stats.json
\newcommand{\NodeTwentySixGlibcProbeThrowsCiLo}{16.38}% data/runs/2026-09-17T08-57-13Z-keypath-runtime-matrix/keypath_stats.json
\newcommand{\NodeTwentySixGlibcVeight}{14}% data/runs/runtime_identity.csv
\newcommand{\NodeTwentySixGlibcVeightFull}{14.6.202.34-node.26}% data/runs/runtime_identity.csv
\newcommand{\NodeTwentySixMuslDecodeOnly}{1.71}% data/runs/2026-09-17T08-57-13Z-keypath-runtime-matrix/keypath_stats.json
\newcommand{\NodeTwentySixMuslHmacString}{3.25}% data/runs/2026-09-17T08-57-13Z-keypath-runtime-matrix/keypath_stats.json
\newcommand{\NodeTwentySixMuslImageId}{a4fb14143ee2}% data/runs/runtime_identity.csv
\newcommand{\NodeTwentySixMuslJwtPreparsed}{6.56}% data/runs/2026-09-17T08-57-13Z-keypath-runtime-matrix/keypath_stats.json
\newcommand{\NodeTwentySixMuslJwtString}{36.42}% data/runs/2026-09-17T08-57-13Z-keypath-runtime-matrix/keypath_stats.json
\newcommand{\NodeTwentySixMuslNodeLiteral}{26.6.0}% data/runs/2026-09-17T08-57-13Z-keypath-runtime-matrix/keypath_stats.json
\newcommand{\NodeTwentySixMuslNodeMajor}{26}% data/runs/2026-09-17T08-57-13Z-keypath-runtime-matrix/keypath_stats.json
\newcommand{\NodeTwentySixMuslOpensslLiteral}{3.5.7}% data/runs/2026-09-17T08-57-13Z-keypath-runtime-matrix/keypath_stats.json
\newcommand{\NodeTwentySixMuslOpensslSeries}{3.5}% data/runs/2026-09-17T08-57-13Z-keypath-runtime-matrix/keypath_stats.json
\newcommand{\NodeTwentySixMuslProbeSucceeds}{15.54}% data/runs/2026-09-17T08-57-13Z-keypath-runtime-matrix/keypath_stats.json
\newcommand{\NodeTwentySixMuslProbeThrows}{22.94}% data/runs/2026-09-17T08-57-13Z-keypath-runtime-matrix/keypath_stats.json
\newcommand{\NodeTwentySixMuslProbeThrowsCiHi}{23.29}% data/runs/2026-09-17T08-57-13Z-keypath-runtime-matrix/keypath_stats.json
\newcommand{\NodeTwentySixMuslProbeThrowsCiLo}{22.79}% data/runs/2026-09-17T08-57-13Z-keypath-runtime-matrix/keypath_stats.json
\newcommand{\NodeTwentySixMuslVeight}{14}% data/runs/runtime_identity.csv
\newcommand{\NodeTwentySixMuslVeightFull}{14.6.202.34-node.26}% data/runs/runtime_identity.csv
\newcommand{\NodeTwentyVeight}{11}% data/runs/runtime_identity.csv
\newcommand{\NodeTwentyVeightFull}{11.3.244.8-node.38}% data/runs/runtime_identity.csv
\newcommand{\PopBestCellPct}{0.0000}% survey/data/population_counts.csv
\newcommand{\PopCapturedAt}{2026-09-18}% survey/data/population_counts.csv
\newcommand{\PopJsImportFiles}{524288}% survey/data/population_counts.csv
\newcommand{\PopJsImportKeyObject}{38}% survey/data/population_counts.csv
\newcommand{\PopJsImportPct}{0.0072}% survey/data/population_counts.csv
\newcommand{\PopJsRequireFiles}{790528}% survey/data/population_counts.csv
\newcommand{\PopJsRequireKeyObject}{29}% survey/data/population_counts.csv
\newcommand{\PopJsRequirePct}{0.0037}% survey/data/population_counts.csv
\newcommand{\PopLargestCellFiles}{790528}% survey/data/population_counts.csv
\newcommand{\PopSumFiles}{1656928}% survey/data/population_counts.csv
\newcommand{\PopSumKeyObject}{168}% survey/data/population_counts.csv
\newcommand{\PopSumOneIn}{9863}% survey/data/population_counts.csv
\newcommand{\PopSumPct}{0.0101}% survey/data/population_counts.csv
\newcommand{\PopTsImportFiles}{333824}% survey/data/population_counts.csv
\newcommand{\PopTsImportKeyObject}{101}% survey/data/population_counts.csv
\newcommand{\PopTsImportPct}{0.0303}% survey/data/population_counts.csv
\newcommand{\PopTsRequireFiles}{8288}% survey/data/population_counts.csv
\newcommand{\PopTsRequireKeyObject}{0}% survey/data/population_counts.csv
\newcommand{\PopTsRequirePct}{0.0000}% survey/data/population_counts.csv
\newcommand{\PopWorstCellPct}{0.0303}% survey/data/population_counts.csv
\newcommand{\ProfEighteenPct}{95.3}% data/runs/2026-09-18T-profile-matched/cpuprof_matched.csv
\newcommand{\ProfEighteenPctHi}{95.6}% data/runs/2026-09-18T-profile-matched/cpuprof_matched.csv
\newcommand{\ProfEighteenPctLo}{95.1}% data/runs/2026-09-18T-profile-matched/cpuprof_matched.csv
\newcommand{\ProfHostPct}{66.2}% data/runs/2026-09-18T-profile-matched/cpuprof_matched.csv
\newcommand{\ProfHostPctHi}{67.4}% data/runs/2026-09-18T-profile-matched/cpuprof_matched.csv
\newcommand{\ProfHostPctLo}{65.3}% data/runs/2026-09-18T-profile-matched/cpuprof_matched.csv
\newcommand{\ProfIterations}{50000}% data/runs/2026-09-18T-profile-matched/cpuprof_matched.csv
\newcommand{\ProfRepeats}{3}% data/runs/2026-09-18T-profile-matched/cpuprof_matched.csv
\newcommand{\ProfTwentySixPct}{61.5}% data/runs/2026-09-18T-profile-matched/cpuprof_matched.csv
\newcommand{\ProfTwentySixPctHi}{61.6}% data/runs/2026-09-18T-profile-matched/cpuprof_matched.csv
\newcommand{\ProfTwentySixPctLo}{58.5}% data/runs/2026-09-18T-profile-matched/cpuprof_matched.csv
\newcommand{\RematchAdjudicated}{12}% survey/data/hand_adjudication.csv
\newcommand{\RematchKeyObject}{0}% survey/data/hand_adjudication.csv
\newcommand{\SampleBuffer}{2}% survey/data/callsites.csv
\newcommand{\SampleCallSites}{327}% survey/data/callsites.csv
\newcommand{\SampleEnvString}{176}% survey/data/callsites.csv
\newcommand{\SampleFileContents}{4}% survey/data/callsites.csv
\newcommand{\SampleKeyObject}{0}% survey/data/callsites.csv
\newcommand{\SampleNoAlgorithms}{309}% survey/data/callsites.csv
\newcommand{\SampleNoAlgorithmsPct}{94}% survey/data/callsites.csv
\newcommand{\SampleRematchAdded}{12}% survey/data/callsites.csv
\newcommand{\SampleRepos}{247}% survey/data/callsites.csv
\newcommand{\SampleStringLiteral}{52}% survey/data/callsites.csv
\newcommand{\SampleUnknown}{93}% survey/data/callsites.csv
\newcommand{\SpanDecodeOnly}{1.69}% data/runs/2026-09-17T08-57-13Z-keypath-runtime-matrix/keypath_stats.json
\newcommand{\SpanHmacString}{1.47}% data/runs/2026-09-17T08-57-13Z-keypath-runtime-matrix/keypath_stats.json
\newcommand{\SpanJwtPreparsed}{1.26}% data/runs/2026-09-17T08-57-13Z-keypath-runtime-matrix/keypath_stats.json
\newcommand{\SpanProbeSucceeds}{11.74}% data/runs/2026-09-17T08-57-13Z-keypath-runtime-matrix/keypath_stats.json
\newcommand{\SpanProbeThrows}{25.36}% data/runs/2026-09-17T08-57-13Z-keypath-runtime-matrix/keypath_stats.json
\newcommand{\SuiteNaiveFailing}{2}% data/runs/2026-09-18T-upstream-suite/suite_results.csv
\newcommand{\SuiteNaivePassing}{509}% data/runs/2026-09-18T-upstream-suite/suite_results.csv
\newcommand{\SuitePatchedFailing}{0}% data/runs/2026-09-18T-upstream-suite/suite_results.csv
\newcommand{\SuitePatchedPassing}{511}% data/runs/2026-09-18T-upstream-suite/suite_results.csv
\newcommand{\SuiteStockFailing}{0}% data/runs/2026-09-18T-upstream-suite/suite_results.csv
\newcommand{\SuiteStockPassing}{511}% data/runs/2026-09-18T-upstream-suite/suite_results.csv

% --- declared but unmeasured ------------------------------------------
% xEightSixReplication: x86_64 replication of the runtime matrix. Not performed: the Oracle Cloud instance was unreachable. See threats to validity.
\newcommand{\xEightSixReplication}{\PLACEHOLDER{xEightSixReplication}}


\usepackage{microtype}
\emergencystretch=3em
\hyphenation{
  micro-ser-vices micro-ser-vice micro-bench-mark micro-bench-marks
  cryp-to-graph-ic cryp-tog-ra-phy au-then-ti-ca-tion va-li-da-tion
  in-stru-men-ta-tion con-fig-u-ra-tion im-ple-men-ta-tion
  Ku-ber-netes side-car side-cars name-space through-put
  pre-parsed re-pro-duc-i-ble de-ploy-ment mea-sure-ment
}
\hypersetup{bookmarksdepth=3}
\raggedbottom

\begin{document}

\title[It Is Not the Cryptography]{It Is Not the Cryptography: A Discarded
Exception Dominates Token Validation Cost in a Node.js Service}

\author[1]{\fnm{Katha} \sur{Sikdar}}\email{ferdouskatha35@gmail.com}
\affil[1]{\orgdiv{Independent Researcher}}

\abstract{The per-request cost of JSON Web Token validation is routinely
attributed to cryptography. We show that in a widely used Node.js library the
attribution is wrong, and that the true mechanism is neither cryptography nor,
as we ourselves previously reported, the conversion of key material. When the
caller supplies a shared secret as a string --- which is what the library's
interface invites --- \texttt{jsonwebtoken} resolves it by attempting to parse it
as an \emph{asymmetric} key first, and falling back to the symmetric constructor
only after that attempt throws. The conversion it falls back to costs
\HostCreateSecretKey~$\mu$s. The discarded failure that precedes it costs
\HostProbeThrows~$\mu$s, or \HostProbeShareOfCall\% of the entire validation call.
Measuring across four container runtimes on one machine, we find the magnitude of
this cost is governed by the bundled OpenSSL version rather than by
containerisation, virtualisation or the JavaScript engine: two runtimes sharing
the OpenSSL \NodeEighteenOpensslSeries{} series but differing in V8 major
version both pay roughly
\NodeEighteenProbeThrows--\NodeTwentyProbeThrows~$\mu$s, while an OpenSSL
\NodeTwentySixMuslOpensslSeries{} runtime pays \NodeTwentySixMuslProbeThrows~$\mu$s. Everything else in the
validation path varies by at most \SpanJwtPreparsed$\times$ across the same
runtimes; the discarded attempt varies by \SpanProbeThrows$\times$. A survey of
public code finds the avoiding pattern essentially absent: of \PopTotalFiles{}
files importing the library, \PopTotalKeyObject{} also mention the pre-parsing
API, and in a sample of \SampleCallSites{} call sites across \SampleRepos{}
repositories, \SampleKeyObject{} pass a pre-parsed key. We propose a fix, and
report a constraint on it that we believe is the more important result: the
discarded attempt is load-bearing for an existing key-type defence, so the
obvious form of the optimisation must not be applied.}

\keywords{JSON Web Token, performance measurement, benchmarking methodology,
Node.js, OpenSSL, software defect}

\maketitle

\section{Introduction}\label{sec:intro}

Authentication in cloud-native services is commonly described as imposing a
per-request cryptographic cost. The reasoning is intuitive: validating a JSON
Web Token means verifying a signature, signatures are cryptography, and
cryptography is expensive relative to parsing. The conclusion usually drawn is
that the cost should be \emph{moved} --- to a sidecar proxy, a gateway, or
dedicated hardware --- because specialised components do cryptographic work well.

We report that for the most widely deployed JavaScript implementation the premise
is wrong, and that our own earlier attempt to correct it was also wrong, in a way
worth recording. That earlier work established that the signature primitive is a
negligible share of the cost and attributed the remainder to converting key
material from a string into a key object on every call. Measuring the conversion
directly shows it costs \HostCreateSecretKey~$\mu$s --- too small to be the
explanation for anything.

The actual mechanism is a control-flow accident. Given a secret that is not
already a key object, the library attempts \texttt{createPublicKey()} first and
falls back to \texttt{createSecretKey()} in the exception handler. For an HMAC
secret the first call cannot succeed. Every validation therefore constructs and
discards an OpenSSL parse failure, at \HostProbeThrows~$\mu$s against a total call
cost of \HostJwtString~$\mu$s.

\paragraph{Contributions}
\begin{enumerate}
\item A decomposition, by two independent methods, locating the dominant cost of
      token validation in a discarded exception rather than in cryptography or
      in key conversion (Section~\ref{sec:mechanism}).
\item Evidence that the magnitude of this cost is governed by the runtime's
      bundled OpenSSL version rather than by containerisation, virtualisation or
      the JavaScript engine, using a runtime pair that separates OpenSSL from V8
      (Section~\ref{sec:runtime}).
\item A survey of public code establishing that the pattern which avoids the cost
      is close to absent in practice (Section~\ref{sec:prevalence}).
\item A fix, together with a constraint on it: the discarded attempt is
      load-bearing for an existing key-type defence, and the natural form of the
      optimisation is therefore unsafe (Section~\ref{sec:fix}).
\item A correction of our own previously reported mechanism, stated explicitly
      rather than quietly revised (Section~\ref{sec:correction}).
\end{enumerate}

\section{Methodology}\label{sec:method}

\subsection{What is measured, and at what granularity}

The unit of measurement is one operating-system process, not one call. Each
process performs \HostCallsPerInvocation{} timed iterations of a single condition
after a warmup of half that number, and contributes its median. Reported figures
are medians of per-process medians, with percentile bootstrap intervals taken
over processes.

This matters because the alternative is misleading in a specific way. An interval
computed over the calls inside one process describes that process's steady state
and is far too narrow: it cannot observe the compilation decisions, inline-cache
states and heap layouts that differ between runs of the same code, which is
precisely what makes a microbenchmark hard to reproduce. Every condition here was
run in \HostInvocations{} separate processes. Between-process coefficient of
variation did not exceed \HostCvMax\% on any condition.

Conditions are interleaved within each round rather than run in blocks. Under
block ordering, a thermal excursion or a background process lands entirely on
whichever condition was executing and is indistinguishable from an effect;
interleaved, it is distributed across conditions and appears as between-round
variance, which the intervals then report.

The clock's own cost is measured in the same process as the condition it
accompanies and is reported rather than subtracted: \HostTimerOverhead~$\mu$s
against the smallest quantity we report, \HostCreateSecretKey~$\mu$s.

\subsection{Decomposing the call}

\texttt{jwt.verify()} is decomposed into the operations it performs, each timed
in isolation under the same protocol: structural decode; the HMAC primitive with
a string key and with a pre-parsed key; the symmetric key constructor; and the
asymmetric key parse that the library attempts first, in both the case where it
throws and the case where it succeeds.

\subsection{Separating the runtime from the deployment}

Our earlier work compared a host microbenchmark against an in-container
measurement and attributed the difference to deployment. Those two environments
differed in more than deployment, so the comparison could not support the
attribution. Here containerisation is held constant and the runtime is varied:
four container images spanning two OpenSSL series and three V8 major versions,
all on one machine, under one container engine.

The decisive pair is Node \NodeEighteenNodeMajor{} and Node \NodeTwentyNodeMajor.
They ship different V8 major versions and share the OpenSSL
\NodeEighteenOpensslSeries{} series. In stock images Node and OpenSSL versions
otherwise co-vary, which would leave the attribution underdetermined.

\subsection{Independent confirmation}

The decomposition above reaches its conclusion by subtracting conditions --- the
same inference style this paper criticises elsewhere. It is therefore checked
against V8's sampling profiler, which attributes self time to named frames with
no arithmetic on our part.

\subsection{Prevalence}

Three measurements with different failure modes. Population counts from code
search over public repositories; a sample of call sites classified by tracing the
key argument within its file; and a targeted counter-search for the avoiding
pattern, fetched in full and adjudicated by hand. The third exists because a
sample can only show a pattern is rare if it could have found it.

\section{The mechanism}\label{sec:mechanism}

\subsection{Where the time goes}

Table~\ref{tab:host} gives the decomposition on the host.

\begin{table}[t]
\caption{Cost of each operation in the validation path. Medians of
per-process medians over \HostInvocations{} processes,
\HostCallsPerInvocation{} calls each; 95\% percentile bootstrap intervals.}
\label{tab:host}
\begin{tabular}{@{}lrr@{}}
\toprule
Operation & Median ($\mu$s) & 95\% CI \\
\midrule
\texttt{verify()}, string secret   & \HostJwtString      & [\HostJwtStringCiLo, \HostJwtStringCiHi] \\
\textbf{failed asymmetric parse}   & \textbf{\HostProbeThrows} & [\HostProbeThrowsCiLo, \HostProbeThrowsCiHi] \\
successful asymmetric parse        & \HostProbeSucceeds  & [\HostProbeSucceedsCiLo, \HostProbeSucceedsCiHi] \\
\texttt{verify()}, pre-parsed key  & \HostJwtPreparsed   & [\HostJwtPreparsedCiLo, \HostJwtPreparsedCiHi] \\
HMAC, string key                   & \HostHmacString     & [\HostHmacStringCiLo, \HostHmacStringCiHi] \\
HMAC, pre-parsed key               & \HostHmacKeyObject  & [\HostHmacKeyObjectCiLo, \HostHmacKeyObjectCiHi] \\
structural decode                  & \HostDecodeOnly     & [\HostDecodeOnlyCiLo, \HostDecodeOnlyCiHi] \\
\textbf{symmetric key constructor} & \textbf{\HostCreateSecretKey} & [\HostCreateSecretKeyCiLo, \HostCreateSecretKeyCiHi] \\
timer overhead                     & \HostTimerOverhead  & [\HostTimerOverheadCiLo, \HostTimerOverheadCiHi] \\
\bottomrule
\end{tabular}
\end{table}

The two bold rows are the result. The conversion to which the cost has been
attributed --- including by us --- costs \HostCreateSecretKey~$\mu$s. The failure
that precedes it costs \HostProbeThrows~$\mu$s, which is
\HostProbeShareOfCall\% of the whole call.

Passing a pre-parsed key instead of a string reduces the call by
\HostStringPenalty~$\mu$s [\HostStringPenaltyCiLo, \HostStringPenaltyCiHi] over
\HostStringPenaltyRounds{} paired rounds.

\subsection{The parts account for most, but not all, of the penalty}

If the penalty were exactly the discarded failure plus the conversion, the two
would sum to it. They do not, quite:

\begin{center}
\begin{tabular}{@{}lr@{}}
\toprule
observed penalty            & \DecompObserved~$\mu$s \\
failed parse $+$ conversion & \DecompPredicted~$\mu$s \\
residual                    & \DecompResidual~$\mu$s \\
\bottomrule
\end{tabular}
\end{center}

The residual's interval, [\DecompResidualCiLo, \DecompResidualCiHi], excludes
zero, so the two measured parts are not the whole story: they account for
\DecompExplainedPct\%, and roughly \DecompResidual~$\mu$s is the library's own
type and claim checking along the non-pre-parsed path. We report this rather than
absorb it into the headline.

\subsection{Confirmation without differencing}

V8's sampling profiler attributes \ProfHostPct\% of self time on the host to the
asymmetric key parse when a string secret is supplied. Under the Node \NodeEighteenNodeMajor{} runtime
the same frame accounts for \ProfEighteenPct\%; under Node \NodeTwentySixMuslNodeMajor, \ProfTwentySixPct\%.
The frame does not appear among the top frames of the pre-parsed path at all. Two
methods, one answer.

\section{The magnitude is a property of the runtime}\label{sec:runtime}

Table~\ref{tab:matrix} varies the runtime with containerisation held constant.

\begin{table}[t]
\caption{The same measurement across four container runtimes on one machine,
\MatrixInvocations{} processes each, \MatrixCallsPerInvocation{} calls per
process. The host row is from Table~\ref{tab:host}.}
\label{tab:matrix}
\begin{tabular}{@{}llrr@{}}
\toprule
Runtime & OpenSSL & Failed parse & Pre-parsed \\
        &         & ($\mu$s)     & \texttt{verify()} \\
\midrule
Node \NodeEighteenNodeMajor, musl  & \NodeEighteenOpensslLiteral     & \NodeEighteenProbeThrows      & \NodeEighteenJwtPreparsed \\
Node \NodeTwentyNodeMajor, musl  & \NodeTwentyOpensslLiteral       & \NodeTwentyProbeThrows        & \NodeTwentyJwtPreparsed \\
Node \NodeTwentySixMuslNodeMajor, musl  & \NodeTwentySixMuslOpensslLiteral  & \NodeTwentySixMuslProbeThrows  & \NodeTwentySixMuslJwtPreparsed \\
Node \NodeTwentySixGlibcNodeMajor, glibc & \NodeTwentySixGlibcOpensslLiteral & \NodeTwentySixGlibcProbeThrows & \NodeTwentySixGlibcJwtPreparsed \\
host (no container) & --- & \HostProbeThrows & \HostJwtPreparsed \\
\bottomrule
\end{tabular}
\end{table}

Three things follow.

\paragraph{It is not containerisation.} Containerised Node \NodeTwentySixGlibcNodeMajor{} on glibc costs
\NodeTwentySixGlibcProbeThrows~$\mu$s against \HostProbeThrows~$\mu$s on the
uncontainerised host, a ratio of \ContainerVersusHostRatio. The container is
marginally \emph{faster}. Whatever produces the large costs in this table, the
act of running inside a container is not it.

\paragraph{It is OpenSSL, not the JavaScript engine.} Node \NodeEighteenNodeMajor{} and Node \NodeTwentyNodeMajor{} ship
different V8 major versions and share the OpenSSL \NodeEighteenOpensslSeries{} series. Both pay of order
\NodeEighteenProbeThrows--\NodeTwentyProbeThrows~$\mu$s. Node \NodeTwentySixMuslNodeMajor{} ships OpenSSL
\NodeTwentySixMuslOpensslSeries{} and pays \NodeTwentySixMuslProbeThrows~$\mu$s on the same libc. Because Node
and OpenSSL versions co-vary in stock images, this attribution is normally
underdetermined; the Node \NodeTwentyNodeMajor{} runtime is what separates them.

\paragraph{The effect is confined to key parsing.} Across every environment
measured, pre-parsed validation spans \SpanJwtPreparsed$\times$, the HMAC
primitive \SpanHmacString$\times$ and structural decode \SpanDecodeOnly$\times$.
The failed parse spans \SpanProbeThrows$\times$. The successful parse moves with
it (\SpanProbeSucceeds$\times$), so this is the cost of OpenSSL \NodeEighteenOpensslSeries{} key parsing
in general rather than anything peculiar to the failure path. The failure is
simply the case in which an application pays that cost on every request for a
result it discards.

\section{How common is the pattern?}\label{sec:prevalence}

A defect that nobody triggers is a curiosity. Three measurements.

\paragraph{Population counts.} Of \PopJsFiles{} public JavaScript files importing
the library, \PopJsKeyObject{} also mention the pre-parsing API; for TypeScript
the figures are \PopTsFiles{} and \PopTsKeyObject{}. Together \PopTotalKeyObject{}
of \PopTotalFiles{}, or \PopKeyObjectPct\% --- about one file in \PopOneIn.

\paragraph{Sampled call sites.} \SampleCallSites{} call sites across
\SampleRepos{} distinct repositories, one file per repository so that no single
project dominates. Environment variables account for \SampleEnvString{}, untraced
identifiers and member expressions \SampleUnknown{}, string literals
\SampleStringLiteral{}, file contents \SampleFileContents{} and buffers
\SampleBuffer{}. Pre-parsed keys account for \SampleKeyObject{}. No file in this
corpus contains the pre-parsing API at all, so the zero reflects the pattern's
absence rather than a classifier that cannot see it.

Incidentally, \SampleNoAlgorithms{} of \SampleCallSites{} call sites
(\SampleNoAlgorithmsPct\%) pass no algorithm restriction. This is not what the
survey set out to measure, but it bears directly on Section~\ref{sec:fix}.

\paragraph{Counter-search, adjudicated by hand.} Every file co-occurring the
library with the pre-parsing API was fetched: \CounterFilesFetched{} files, of
which \CounterWithVerify{} contain a validation call. An automated pass flagged
\AdjudicatedTotal{} as passing a pre-parsed key; all were then read. Of these,
\AdjudicatedFalsePositive{} are false positives --- projects defining their own
helper function of the same name, one of which returns a string --- and
\AdjudicatedUnresolved{} could not be resolved within the file. Genuine:
\AdjudicatedGenuine{} call sites across \AdjudicatedGenuineProjects{} projects.

\section{A fix, and a constraint on it}\label{sec:fix}

The fix is to dispatch on the token's declared algorithm before attempting the
asymmetric parse, so that an HMAC token with an ordinary string secret is
resolved directly. Applied, this takes the call from \HostJwtString~$\mu$s to
\HostJwtStringSafe~$\mu$s, a factor of \HostSafePatchRatio{} and within
\HostSafePatchResidual~$\mu$s of supplying a pre-parsed key, with no change in
calling code.

The constraint is the more important half. The discarded attempt is not purely
wasteful: the \emph{success} of that parse on asymmetric material is what
produces a key object whose type a later check depends on, and that check is what
prevents symmetric and asymmetric key material from being interchanged. A change
that dispatches on the declared algorithm while relaxing which key material may
take the fast path alters which inputs reach that check.

The fix must therefore be narrow: the fast path is restricted to string secrets
that cannot be asymmetric key material, so that anything which might be resolves
exactly as it does today. We verified that the narrow form preserves the check
and that a form without the material restriction does not. Two observations
sharpen the point. First, a regression test that passes on both a correct and an
incorrect build constrains nothing; the test reported here fails on the
unrestricted form, which is what makes it evidence rather than decoration.
Second, the independent defence against this class of confusion is for the caller
to restrict permitted algorithms --- and \SampleNoAlgorithmsPct\% of surveyed call
sites do not.

\section{Correcting our own earlier report}\label{sec:correction}

Our earlier manuscript reported that the dominant term was key material converted
from a string on every call. That is wrong: the conversion costs
\HostCreateSecretKey~$\mu$s. It also explained an asymmetry between HMAC and RSA
costs as conversion being dearer than parsing a PEM, which is backwards --- a
successful asymmetric parse costs \HostProbeSucceeds~$\mu$s against
\HostCreateSecretKey~$\mu$s for the conversion. The real asymmetry is that the
RSA path's parse \emph{succeeds}, making its cost useful work, while the HMAC
path's fails, making its cost waste.

That manuscript further compared a host microbenchmark against an in-container
measurement and attributed the difference to deployment. Section~\ref{sec:runtime}
shows the comparison spanned two OpenSSL major series, so it could not support
that attribution. It also stated a runtime version for the system under test that
was read from the host rather than from the container, and was wrong by two major
versions.

We record these rather than revise them silently because the errors share a
structure worth naming: in each case a quantity was attributed to the variable we
were interested in, when a second variable had changed at the same time and had
not been recorded at the granularity that would have revealed it.

\section{Threats to validity}\label{sec:threats}

\paragraph{One machine, one instruction set.} Every measurement in this paper
comes from a single ARM64 host: an Apple silicon laptop, its container runtimes
under one engine. We have \textbf{not} verified any of this on x86\_64.

We state what we expect and why, and then state that we did not test it. The
mechanism is a property of a library's API contract --- an asymmetric parse
attempted before a symmetric one, with the failure discarded --- and of the cost
of key parsing in a particular OpenSSL series. Neither is a property of an
instruction set, and we know of no reason for the control flow to differ on
x86\_64. But an expectation is not a measurement. The magnitude in particular
depends on how OpenSSL's key parsing performs on a given platform, and that is
exactly the kind of quantity that does vary with architecture. Until the
following is filled in, no claim in this paper should be read as holding beyond
ARM64: \xEightSixReplication.

\paragraph{One library, one version.} The mechanism is specific to one JavaScript
implementation at one version. Whether other implementations or other language
ecosystems attempt key resolution in the same order is unmeasured.

\paragraph{The prevalence survey measures files, not systems.} A project that
constructs a pre-parsed key in one file and validates in another is invisible to
file-level analysis, and this is the strongest threat to
Section~\ref{sec:prevalence}. Within-file tracing found no such case, which is
weak evidence against it. Code search covers public repositories only and
includes forks, vendored dependencies and teaching material; the numerator is
inflated in the same way as the denominator, but the ratio is not a clean
population proportion and should not be read as one. Public code also skews
toward small projects, and the \AdjudicatedGenuineProjects{} projects found using
the API correctly are all small.

\paragraph{Absolute figures are from a shared laptop.} The host also ran a
desktop environment. Contention inflates a duration rather than deflating it, so
the reported costs are upper bounds and the comparisons between conditions ---
which is where every claim here rests --- are unaffected in direction. Absolute
values should not be compared against dedicated hardware.

\paragraph{The host's cryptographic library moved mid-study.} Installing an
unrelated command-line tool upgraded the host's OpenSSL beneath an unchanged Node
binary, because the runtime links it dynamically. No reported measurement is
affected: the affected runs predate the upgrade and recorded their OpenSSL
version, and container measurements bundle their own. It was detectable only
because that version is recorded per measurement rather than per run, which we
note as a methodological recommendation rather than an incident.

\section{Conclusion}\label{sec:conclusion}

The per-request cost of token validation in this library is real, is worth
attending to, and is not cryptographic. It is not key conversion either. It is a
parse attempted in the wrong order, whose failure is constructed and thrown away
on every request, and whose price is set by the OpenSSL version a base image
happens to pin --- varying \SpanProbeThrows$\times$ across runtimes while the rest
of the validation path varies by at most \SpanJwtPreparsed$\times$. The pattern
that avoids it is close to absent from public code.

The methodological lesson is not that one should measure in deployment. It is
narrower and, we think, more useful: attribute a cost only to a variable you
recorded at the granularity at which it can change. Both errors we correct here
--- our own, and the field's --- have that shape.

\section*{Declarations}

\bmhead{Data and code availability} All measurements, the instrumentation, the
analysis and the scripts that produce every numeric value in this paper are in
the replication package. Every number is generated from a file; none is entered
by hand. Quantities that were not measured appear as red placeholders rather than
as absences.

\bmhead{Competing interests} The author declares no competing interests.

\end{document}
